\section{Powiązania z innymi działami matematyki}\label{sec:powiazania}

\subsection{Algebra: od jąder do kategorii abelowych}

Kategoria $\Vect_k$ ma obiekt zerowy, jądra i kojądra wszystkich morfizmów, a każdy zbiór morfizmów $\Hom(V,W)$ jest grupą abelową względem dodawania odwzorowań liniowych. Ponadto składanie jest addytywne w każdej zmiennej:
$(f_1+f_2)\circ g=f_1\circ g+f_2\circ g$ oraz $f\circ(g_1+g_2)=f\circ g_1+f\circ g_2$.

\begin{definicja}[Kategorie addytywne i abelowe]\label{cat:def:abelian-category}
Kategoria preaddytywna to kategoria lokalnie mała, której zbiory morfizmów są grupami abelowymi, a składanie jest addytywne w obu zmiennych. Kategoria addytywna jest kategorią preaddytywną z obiektem zerowym i skończonymi biproduktami (obiektami będącymi jednocześnie iloczynami i koproduktami, z kanonicznymi projekcjami i włączeniami). Dokładniej dla biproduktu $P=A\oplus B$ wymagamy, aby projekcje $p_A,p_B$ i włączenia $i_A,i_B$ spełniały
\[
p_Ai_A=\id_A,\quad p_Bi_B=\id_B,\quad
p_Ai_B=0,\quad p_Bi_A=0,\quad
i_Ap_A+i_Bp_B=\id_P.
\]
Równości te zapewniają obie własności uniwersalne: morfizm do $P$ o składowych $u,v$ to $i_Au+i_Bv$, a morfizm z $P$ o ograniczeniach $r,s$ to $rp_A+sp_B$; jednoznaczność w obu przypadkach wynika z ostatniej równości. Kategoria abelowa jest kategorią addytywną, w której każdy morfizm ma jądro i kojądro, każdy monomorfizm jest jądrem pewnego morfizmu, a każdy epimorfizm jest kojądrem pewnego morfizmu.
\end{definicja}
Jądro $f:A\to B$ jest equalizatorem $f$ i morfizmu zerowego $A\to B$, a kojądro jest ich koequalizatorem. Obraz definiujemy jako jądro kojądra, traktowane jako podobiekt przeciwdziedziny. Ciąg $A\xrightarrow{f}B\xrightarrow{g}C$ w kategorii abelowej jest dokładny w $B$, gdy obraz $f$ i jądro $g$ są tym samym podobiektem $B$ (tj. ich morfizmy do $B$ są zgodne poprzez izomorfizm dziedzin). W $\Vect_k$ i $\Ab$ są to odpowiednio zwykłe jądro, rzut na iloraz przez obraz oraz włączenie obrazu.

\begin{stwierdzenie}\label{cat:prop:abelian-examples}
Kategorie $\Vect_k$ i $\Ab$ są abelowe.
\end{stwierdzenie}
\begin{proof}
W obu kategoriach dodawanie morfizmów jest punktowe i składanie jest dwuliniowe względem grup abelowych morfizmów. Obiekt zerowy to przestrzeń lub grupa trywialna. Dla dwóch obiektów $A,B$ suma prosta $A\oplus B$ jest iloczynem: morfizmy $u:X\to A$, $v:X\to B$ wyznaczają jedyny $x\mapsto(u(x),v(x))$. Jest też koproduktem: morfizmy $r:A\to Y$, $s:B\to Y$ wyznaczają jedyny $(a,b)\mapsto r(a)+s(b)$, ponieważ $(a,b)=(a,0)+(0,b)$. Projekcje i włączenia współrzędnych spełniają pięć podanych równości, co sprawdzamy na dowolnej parze $(a,b)$. Obiekt zerowy obsługuje pusty biprodukt, a iteracja konstrukcji daje wszystkie skończone biprodukty.

Jądro $f$ jest włączeniem $\Ker f\to A$: morfizm $h$ spełniający $fh=0$ ma wartości w $\Ker f$, więc faktoryzuje się jednoznacznie przez to włączenie. Kojądro jest rzutem $q:B\to B/\im f$: dla $g:B\to Z$ z $gf=0$ wzór $\bar g(b+\im f)=g(b)$ jest dobrze określony, bo $g$ zeruje się na $\im f$; daje jedyną faktoryzację przez surjekcję $q$.

Monomorfizmy w obu kategoriach są dokładnie morfizmami injektywnymi: injekcja daje skracanie z lewej, a przy niezerowym elemencie jądra dwa różne morfizmy z $k$ (odpowiednio $\mathbb Z$), wysyłające generator na ten element i na zero, mają to samo złożenie. Epimorfizmy są dokładnie surjekcjami: surjekcja daje skracanie z prawej; gdy obraz jest właściwy, rzut na niezerowy iloraz przez obraz i morfizm zerowy są różne, lecz ich złożenia z $f$ są równe. Monomorfizm identyfikuje dziedzinę izomorficznie z obrazem, więc jest jądrem rzutu na iloraz przez obraz. Dla epimorfizmu $f$ odwzorowanie $A/\Ker f\to B$, $a+\Ker f\mapsto f(a)$, jest dobrze określoną bijekcją liniową lub izomorfizmem grup: injektywność wynika z definicji jądra, surjektywność z surjektywności $f$. Zatem $f$ jest kojądrem włączenia jądra. Spełniono wszystkie warunki definicji.
\end{proof}

\begin{przyklad}[Ciąg dokładny związany z morfizmem]\label{cat:ex:exact-sequence}
Dla liniowego $f:V\to W$ mamy ciąg $0\to\Ker f\xrightarrow{i}V\xrightarrow{f}W\xrightarrow{q}W/\im f\to0$. Jest dokładny: $\Ker i=0$, $\im i=\Ker f$, $\Ker q=\im f$, a $q$ jest surjektywny. W tej postaci nie jest to krótki ciąg dokładny (krótki ciąg ma trzy obiekty między skrajnymi zerami). Dwa związane z nim krótkie ciągi dokładne to
\[
0\to\Ker f\to V\to\im f\to0,\qquad
0\to\im f\to W\to W/\im f\to0.
\]
W pierwszym morfizm $V\to\im f$ jest zawężeniem przeciwdziedziny $f$, w drugim $\im f\to W$ jest włączeniem. Dokładność wynika z tych samych równości jąder i obrazów. Por. skrypt \emph{Topologia algebraiczna}, rozdział o algebrze homologicznej.
\end{przyklad}

\subsection{Topologia: funktorialność niezmienników}

Z punktu widzenia teorii kategorii wielu niezmienników topologicznych nie należy traktować jako pojedynczych liczb lub grup, lecz jako funktory. Przykładowo grupa podstawowa przyporządkowuje \emph{punktowanej} przestrzeni topologicznej $(X,x_0)$ grupę $\pi_1(X,x_0)$. Każde ciągłe odwzorowanie punktowane $f:(X,x_0)\to(Y,y_0)$ indukuje homomorfizm $\pi_1(f):\pi_1(X,x_0)\to\pi_1(Y,y_0)$ przez $[\gamma]\mapsto[f\circ\gamma]$. Jeżeli pętle są homotopijne przy ustalonych końcach, złożenie homotopii z $f$ zachowuje tę własność, więc przyporządkowanie klas jest dobrze określone. Złożenie $f$ z konkatenacją pętli jest konkatenacją ich obrazów, stąd homomorfizm. Identyczności i składanie zachowuje, ponieważ $g\circ(f\circ\gamma)=(g\circ f)\circ\gamma$. Jest więc funktorem z kategorii punktowanych przestrzeni topologicznych do $\Grp$.

\begin{przyklad}[Homologia jako funktor]
Dla każdego $n\geq0$ homologia singularna o współczynnikach całkowitych jest funktorem $H_n:\Top\to\Ab$. Na wolnych grupach abelowych $C_n(X)$ generowanych przez sympleksy $\sigma:\Delta^n\to X$ definiujemy $C_n(f)(\sigma)=f\circ\sigma$ i rozszerzamy addytywnie. Dla $n\geq1$ i włączeń ścian $\delta_i:\Delta^{n-1}\to\Delta^n$ mamy
\[
\partial C_n(f)(\sigma)=\sum_i(-1)^i f\circ\sigma\circ\delta_i
=C_{n-1}(f)(\partial\sigma).
\]
Dla $n=0$ różniczka jest zerowa. Zatem cykle przechodzą na cykle, a brzegi na brzegi, więc wzór $H_n(f)([c])=[C_n(f)(c)]$ jest dobrze określony: zmiana reprezentanta o brzeg zmienia obraz o brzeg. Na generatorach zachodzą $C_n(\id_X)=\id_{C_n(X)}$ i $C_n(g\circ f)=C_n(g)C_n(f)$, zatem te same równości zachodzą na ilorazach homologii. Por. skrypt \emph{Topologia algebraiczna}, rozdział o homologii przestrzeni. To dlatego homeomorficzne przestrzenie mają izomorficzne grupy homologii: stosujemy ogólne twierdzenie~\ref{tw:funktory-iso} o zachowywaniu izomorfizmów przez funktory.
\end{przyklad}

\subsection{Analiza: przestrzenie funkcji jako funktory Hom}

Jeśli rozważamy kategorię odpowiednio dobranych przestrzeni analitycznych i dozwolonych odwzorowań, to zapis $\Hom(X,Y)$ oznacza \emph{przestrzeń odwzorowań zgodnych ze strukturą}, np. ciągłych lub holomorficznych. Własności uniwersalne mogą wtedy opisywać obiekty nie przez ich punkty, lecz przez wszystkie rodziny odwzorowań do nich i z nich. Zapis $\Hom(X,Y)$ sam w sobie oznacza zbiór morfizmów; dodatkowa topologia lub struktura liniowa wymaga osobnego określenia. To jest podstawowy punkt widzenia geometrii algebraicznej, teorii schematów i geometrii analitycznej.

\begin{przyklad}[Iloraz przestrzeni unormowanej jako konstrukcja uniwersalna]\label{cat:ex:normed-quotient}
Niech $X$ będzie przestrzenią unormowaną nad $k=\mathbb R$ lub $\mathbb C$, a $M$ jej domkniętą podprzestrzenią liniową. W kategorii przestrzeni unormowanych i ograniczonych odwzorowań liniowych rzut $q:X\to X/M$ z normą ilorazową $\|x+M\|=\inf_{m\in M}\|x+m\|$ jest kojądrem włączenia $M\to X$ (względem morfizmu zerowego). Istotnie, dla ograniczonego $f:X\to Y$ zerującego się na $M$ wzór $\bar f(x+M)=f(x)$ jest dobrze określony i liniowy. Dla każdego $m\in M$ mamy $\|\bar f(x+M)\|=\|f(x+m)\|\leq\|f\|\|x+m\|$, więc po wzięciu infimum $\|\bar f\|\leq\|f\|$. Ponieważ $f=\bar f q$ i $\|q\|\leq1$, również $\|f\|\leq\|\bar f\|$. Jednoznaczność wynika z surjektywności $q$. Por. skrypt \emph{Analiza funkcjonalna}, rozdział o konstrukcjach ilorazowych.
\end{przyklad}

\begin{przyklad}[Kategoria przestrzeni unormowanych nie jest abelowa]\label{cat:ex:normed-nonabelian}
Niech $c_{00}$ oznacza przestrzeń ciągów o skończenie wielu niezerowych wyrazach, z normą $\ell^2$. Włączenie $j:c_{00}\to\ell^2$ jest ograniczone i injektywne, więc jest monomorfizmem. Jest także epimorfizmem: jeśli ograniczone odwzorowania liniowe $g,h:\ell^2\to Y$ spełniają $gj=hj$, to dla każdego $x\in\ell^2$ jego obcięcia $x^{(N)}\in c_{00}$ spełniają
\[
\|x-x^{(N)}\|_2^2=\sum_{n>N}|x_n|^2\longrightarrow0.
\]
Z ciągłości $g,h$ otrzymujemy $g(x)=\lim_N g(x^{(N)})=\lim_N h(x^{(N)})=h(x)$. Jednak $j$ nie jest surjekcją: ciąg $(2^{-n})_{n\geq1}$ należy do $\ell^2$ i nie należy do $c_{00}$. Nie jest zatem izomorfizmem.

Gdyby ta kategoria była abelowa, $j$ jako monomorfizm byłby jądrem pewnego morfizmu $r$. Mielibyśmy $rj=0$, więc z epiczności $j$ wynikałoby $r=0$. Jądro morfizmu zerowego o dziedzinie $\ell^2$ jest identycznością $\ell^2$; z jednoznaczności jąder z dokładnością do izomorfizmu wynikałoby, że $j$ jest izomorfizmem. Sprzeczność. Przykład pokazuje, dlaczego algebraiczne własności jąder i kojąder nie przenoszą się automatycznie na kategorie z warunkiem ciągłości.
\end{przyklad}

\subsection{Geometria: presnopy i snopy}

\begin{przyklad}[Presnop funkcji ciągłych]\label{cat:ex:continuous-sheaf}
Dla przestrzeni topologicznej $X$ kategoria $\mathcal O(X)$ ma zbiory otwarte jako obiekty i inkluzje jako morfizmy. Presnop zbiorów jest funktorem $\mathcal O(X)^{\op}\to\Set$. Przykładem jest $\mathcal F(U)=C(U,\mathbb R)$, z ograniczeniami $r_{U,V}:\mathcal F(U)\to\mathcal F(V)$ dla $V\subseteq U$. Zachodzą $r_{U,U}=\id$ i $r_{V,W}r_{U,V}=r_{U,W}$ dla $W\subseteq V\subseteq U$.

Presnop jest snopem, jeśli dla każdego pokrycia otwartego $U=\bigcup_i U_i$ rodzina $s_i\in\mathcal F(U_i)$ zgodna po ograniczeniu do $U_i\cap U_j$ ma dokładnie jedno sklejenie w $\mathcal F(U)$. Dla funkcji ciągłych definiujemy $s(x)=s_i(x)$, gdy $x\in U_i$; zgodność daje niezależność od wyboru $i$. Dla otwartego $O\subseteq\mathbb R$ mamy $s^{-1}(O)=\bigcup_i s_i^{-1}(O)$, otwarty w $U$, bo każdy $U_i$ jest otwarty w $U$. Zatem $s$ jest ciągła. Jednoznaczność wynika z pokrycia, również dla $U=\varnothing$ (istnieje dokładnie jedna funkcja pusta). Tak warunek lokalnego sklejania uzupełnia samą funktorialność ograniczeń. Por. skrypt \emph{Geometria algebraiczna}, rozdział o presnopach i snopach.
\end{przyklad}

\begin{uwaga}[Dalsza ścieżka]
Po opanowaniu podstaw warto kolejno poznać kategorie addytywne i abelowe, kategorie monoidalne, sprzężenia i monady w większej ogólności, algebry nad monadą, a następnie kategorie pochodne lub kategorie wyższe. Nie są one wymagane do zrozumienia definicji funktora, transformacji naturalnej czy lematu Yonedy, lecz rozwijają te same idee.
\end{uwaga}

